\subsection{Ten cycles on six powers and eleven cycles on seven powers in degree eighteen}\label{subsec:six-power-proof}

The balanced trinomial already has two compact brackets and two
incoming tests. The fold of Article, Proposition~4.4 closes
the incoming trajectories in degree eighteen. The second construction
uses the empty interval between the origin and the innermost of these
cycles to add one more power.

\begin{proof}[Proof of the $k=6$ row of Theorem~\ref{thm:lower}\textup{(ii)}]
Take the monic balanced degree-ten seed $P_{\mathrm{seed},\lambda}(x)=x^{10}+bx^9+\lambda x^5$
($d_{\mathrm{top}}=10$, $m_0=5$) from Article, Lemma~3.1(ii)--(iii). It provides
$N_{\mathrm{ord}}=2$ bounded ordinary sign brackets and two negative heights
$0>h_{\mathrm{seed},1}>h_{\mathrm{seed},2}$ at which $G_+(h)-G_-(h)$ at $\lambda_0$ has signs $-,+$.
Fix $s>0$ so small that $G_+(h_{\mathrm{seed},1})-G_-(h_{\mathrm{seed},1}+s)<0$ and
$G_+(h_{\mathrm{seed},2})-G_-(h_{\mathrm{seed},2}+s)>0$, and choose
$h_{\mathrm{deep}}<h_{\mathrm{seed},2}$ so negative that $|G_\pm(h)-h|<s/4$ for all $h\le
h_{\mathrm{deep}}+s$ (by Article, Lemma~2.2(iv)), ensuring
$G_+(h_{\mathrm{deep}})-G_-(h_{\mathrm{deep}}+s)<-s/2<0$. Apply Article, Proposition~4.4 to
$(h_1,h_2,h_3)=(h_{\mathrm{seed},1},h_{\mathrm{seed},2},h_{\mathrm{deep}})$ with signs $(-,+,-)$.
From $\beta_{\rm refl}$ to $\beta_{\rm ext}$ the signs $(+,-,+,-,+)$
give four disjoint brackets. With the two bounded seed brackets, two
fold brackets and two marked cycles, Article, Proposition~4.4
gives ten simultaneous hyperbolic limit cycles on $\{5,9,10\}\cup\{11,13,18\}=\{5,9,10,11,13,18\}$.
\end{proof}

\begin{corollary}[Explicit degree-eighteen seven-power support]\label{cor:seven-inner}
A nonempty open subset of $(\mathbb{R}\setminus\{0\})^7$ on the degree-$18$ support
$\{3,5,9,10,11,13,18\}$ realizes at least $11$ simultaneous hyperbolic limit cycles.
\end{corollary}

\begin{proof}
Fix a six-term polynomial $F_0$ on $\{5,9,10,11,13,18\}$ with ten hyperbolic limit cycles, and isolate
their positive $y$-axis intercepts in pairwise disjoint compact intervals. Below the innermost
interval, the linearized rotation gives two transverse returns for
small launches. Choose $\beta_0>0$ so both arcs lie in a neighborhood of
$(0,0)$ where the odd part of $F_0$ has the strict sign of its lowest odd coefficient $a_5\ne 0$; by
Article, Lemma~2.1(ii), $\oscD_{F_0}(\beta_0)$ has the strict sign $\sgn(a_5)$. Adjoin
$\gamma x^3$ of opposite sign $\sgn(\gamma)=-\sgn(a_5)$ with $|\gamma|>0$ small enough that all ten
hyperbolic cycles and the strict sign at $\beta_0$ persist. At a still smaller height
$\beta_1\in(0,\beta_0)$, the cubic term $\gamma r^3$ dominates the odd part on both half-arcs, giving
$\sgn\oscD(\beta_1)=\sgn(\gamma)=-\sgn\oscD(\beta_0)$ and bracketing an eleventh periodic orbit in
$(\beta_1,\beta_0)$. Applying Article, Lemma~2.3 to the $x^3$ coefficient makes all eleven
orbits simultaneously hyperbolic while keeping all seven powers in $\{3,5,9,10,11,13,18\}$ nonzero.
\end{proof}
